\chapter*{Literatur und Einordnung}
\addcontentsline{toc}{chapter}{Literatur und Einordnung}
\markboth{LITERATUR UND EINORDNUNG}{LITERATUR UND EINORDNUNG}

Dieses Verzeichnis ist bewusst kommentiert. Die Bausteine der
Horimetrik -- Mischbasensysteme, fallende Fakultäten,
Inversionsfolgen, das Differenzenkalkül -- sind etablierte
Mathematik, und dieses Buch behauptet nichts anderes. Jeder Eintrag
nennt deshalb zweierlei: welchen Baustein er beisteuert, und was
dort jeweils \emph{nicht} steht -- also die Stelle, an der die
Horimetrik über ihre Quelle hinausgeht. Wer die Einordnung aus
Kapitel~\ref{ch:einordnung} nachprüfen will, findet hier die
Belege.

% thebibliography setzt sonst ein eigenes \chapter* ("Literaturverzeichnis")
% auf eine neue Seite; wir haben oben schon eine Kapitelüberschrift.
\begingroup
\makeatletter
\renewcommand{\chapter}[2]{}%
\renewcommand{\@mkboth}[2]{}%
\makeatother
\begin{thebibliography}{10}

\bibitem{cantor} G.~Cantor: \emph{Über die einfachen Zahlensysteme},
Zeitschrift für Mathematik und Physik 14 (1869), 121--128.
--- Die Fakultätsbasis für Brüche; das ist exakt der bruchtreue
Grenzraum der dritten Treppe (Kapitel~\ref{ch:erweiterungen}),
nicht aber die gestalttreue Welt oder das Zusammenspiel.

\bibitem{knuth} D.~E. Knuth: \emph{The Art of Computer Programming},
Band~2, 3.~Aufl., Addison--Wesley 1997, Abschnitt 4.1.
--- Mischbasensysteme und das Fakultätszahlensystem (Factoradic);
dort wachsen die Basen zur signifikanten Seite, führende Nullen sind
wertneutral.

\bibitem{lehmer} D.~H. Lehmer: \emph{Teaching combinatorial tricks
to a computer}, Proc. Sympos. Appl. Math. 10 (1960), 179--193.
--- Lehmer-Codes und Inversionstabellen; erklären die
Fakultätsgröße der Horizonte und die Ziffernschranke $0\le a_i\le i$,
nicht aber die horizontabhängige Auswertung.

\bibitem{gkp} R.~L. Graham, D.~E. Knuth, O.~Patashnik:
\emph{Concrete Mathematics}, 2.~Aufl., Addison--Wesley 1994, Kap.~2.
--- Fallende Fakultäten und Differenzenkalkül; das Werkzeug der
Horizontableitung.

\bibitem{cahen} P.-J. Cahen, J.-L. Chabert: \emph{Integer-Valued
Polynomials}, AMS Surveys and Monographs 48, 1997.
--- Polynome mit ganzzahligen Werten in der Binomialbasis; die
Horimetrik nutzt nur den positiven Kegel und wertet am laufenden
Horizont aus.

\bibitem{rigo} M.~Rigo, M.~Stipulanti: \emph{Revisiting regular
sequences in light of rational base numeration systems},
arXiv:2103.16966.
--- Abstrakte Numerationssysteme, in denen führende Nullen über die
Radix-Ordnung wertsensitiv sind: der nächste bekannte Nachbar der
Nullen-Semantik, dort formalsprachlich statt algebraisch behandelt.

\bibitem{kudr} G.~Kudryavtseva, V.~Mazorchuk: \emph{On Kiselman's
semigroup}, arXiv:math/0511374; O.~Ganyushkin, V.~Mazorchuk:
\emph{On Kiselman quotients of 0-Hecke monoids}, arXiv:1006.0316.
--- Monoide aus idempotenten Erzeugern mit Absorptionsrelationen;
Nachbarschaft des sechselementigen Kompaktifizierungsmonoids.

\bibitem{simon} D.~Simon: \emph{Mixed radix numeration bases:
Horner's rule, Yang-Baxter equation and Furstenberg's conjecture},
arXiv:2405.19798.
--- Lokale Basiswechsel in Mischbasensystemen; Kandidatenwerkzeug
für den Interchange-Defekt.

\bibitem{invseq} S.~Corteel, M.~A.~Martinez, C.~D.~Savage,
M.~Weselcouch: \emph{Patterns in Inversion Sequences I},
arXiv:1510.05434.
--- Inversionsfolgen ($0\le e_j<j$) und ihre Bijektion zu
Permutationen; unser $H_n$ ist als Menge genau dieser Raum
(Kapitel~\ref{ch:einordnung}). Was dort nicht steht: die umgedrehte
Wertung, die Horizontübergänge und der Strukturhorizont als
Statistik.

\bibitem{lecturehall} C.~D. Savage, M.~J. Schuster: \emph{Ehrhart
series of lecture hall polytopes and Eulerian polynomials for
inversion sequences}, J. Combin. Theory Ser.~A 119 (2012),
850--870.
--- $s$-Inversionsfolgen mit beliebigen Kapazitätsprofilen und ihre
Geometrie; der Rahmen für das $s$-Horimetrik-Programm aus
Kapitel~\ref{ch:einordnung}.

\bibitem{tightentries} S.~Chern, S.~Fu, Z.~Lin: \emph{Burstein's
permutation conjecture, Hong and Li's inversion sequence
conjecture, and restricted Eulerian distributions},
arXiv:2209.12137.
--- Statistiken auf Inversionsfolgen, darunter die \emph{tight
entries} ($e_i=s_i-1$, Statistik $\mathrm{tig}$); der
Strukturhorizont verfeinert diese Extremfall-Zählung zu einem
Abstandsmaß (Kapitel~\ref{ch:einordnung}).

\bibitem{middleorder} M.~Bouvel, L.~Ferrari, B.~E.~Tenner:
\emph{Between weak and Bruhat: the middle order on permutations},
arXiv:2405.08943.
--- Die koordinatenweise Ordnung auf Inversionscodes als Gitter
zwischen schwacher und Bruhat-Ordnung; der Rahmen, in dem
$\mathcal F_r$ ein Hauptideal und $\sigma$ eine Idealstufe wird.

\bibitem{umbral} K.~Beauduin: \emph{Operational Umbral Calculus},
arXiv:2407.16348.
--- Moderne finite Operatorrechnung: Delta-Operatoren und ihre
Basisfolgen. Der Rahmen für $\Delta$ und die Strukturpolynome; was
dort nicht steht, ist die Kapazitätsbuchhaltung $\Gamma_T$.

\bibitem{ikenmeyerpak} C.~Ikenmeyer, I.~Pak: \emph{What is in \#P
and what is not?}, FOCS 2022, arXiv:2204.13149.
--- Charakterisiert die \emph{binomial-good}-Polynome
(nichtnegative ganzzahlige Binomialbasis-Koeffizienten) als die
\emph{relativierenden} polynomiellen \#P-Abschlussoperationen und
formuliert die Binomial-Basis-Vermutung (unrelativiert gilt
dieselbe Charakterisierung nur vermutungsweise), deren univariate
Fassung bereits P$\,\ne\,$NP implizieren würde. Die Gestalten der Horimetrik sind wegen
$\ff Xd=d!\binom Xd$ eine echte Teilklasse der binomial-good-Polynome
(Kapitel~\ref{ch:strukturpolynome}); was dort \emph{nicht} steht,
ist die quantitative Schicht: die Höhe $\max(d+c_d)$, die
$(r{+}1)!$-Hierarchien und deren Wachstumsgesetze.

\bibitem{oeis} OEIS Foundation Inc.: \emph{The On-Line Encyclopedia
of Integer Sequences}, \texttt{https://oeis.org}.
--- Referenzdatenbank für die neun Zählfolgen dieses Buchs;
A-Nummern werden nach Vergabe nachgetragen.

\end{thebibliography}
\endgroup
